\documentclass[a4paper, 10pt]{article}

\usepackage[utf8]{inputenc}
\usepackage[T1]{fontenc}
\usepackage{indentfirst} % first line indentation
\usepackage{pifont} % hand-right icon
\usepackage{marginnote} % notes on margins
\usepackage{sectsty}
\usepackage{graphicx}
\usepackage{amssymb}
\usepackage{amsmath}
\usepackage{bm}
\usepackage[english]{babel}
\usepackage{biblatex}
\addbibresource{references.bib}
\usepackage{xcolor}
\usepackage{listings} % insert code snippets
\usepackage{hyperref}

% Graphics path
\graphicspath{{../../../fig/}}

% Margins
\topmargin=-0.45in
\evensidemargin=0in
\oddsidemargin=0in
\textwidth=6.5in
\textheight=9.0in
\headsep=0.25in

% Bibliography style
\title{Image Processing in \textsc{Python}}
\author{Prof. Dr. Aldo André Díaz-Salazar}
\date{}

% User definitions
\newcommand{\lt}{\ensuremath <}
\newcommand{\gt}{\ensuremath >}
\newcommand{\Figureref}[1]{\hyperref[#1]{Figure~\ref*{#1}}}

\hypersetup{
	%pagebackref=true,
	pdftitle={\@title},
	pdfauthor={\@author},
	pdfcreator={LaTeX with abnTeX2},
	pdfkeywords={abnt}{latex}{abntex}{abntex2}{academic work},
	hidelinks,					% desabilita as bordas dos links
	colorlinks=true,       	% false: boxed links; true: colored links
	linkcolor=blue,          	% color of internal links
	citecolor=blue,        		% color of links to bibliography
	filecolor=magenta,      	% color of file links
	urlcolor=blue,
%	linkbordercolor={1 1 1},	% set to white
	bookmarksdepth=4
}

% Space between matrix brackets
\makeatletter
\renewenvironment{bmatrix}
{\left[\mkern10mu\env@matrix}
{\endmatrix\mkern10mu\right]}

\definecolor{codegreen}{rgb}{0,0.6,0}
\definecolor{codegray}{rgb}{0.5,0.5,0.5}
\definecolor{codepurple}{rgb}{0.58,0,0.82}
\definecolor{backcolour}{rgb}{0.95,0.95,0.92}

\lstdefinestyle{mystyle}{
	backgroundcolor=\color{backcolour},   
	commentstyle=\color{codegreen},
	keywordstyle=\color{magenta},
	numberstyle=\tiny\color{codegray},
	stringstyle=\color{codepurple},
	basicstyle=\ttfamily\footnotesize,
	breakatwhitespace=false,         
	breaklines=true,                 
	captionpos=b,                    
	keepspaces=true,                 
	numbers=left,                    
	numbersep=5pt,                  
	showspaces=false,                
	showstringspaces=false,
	showtabs=false,                  
	tabsize=2
}

\lstset{style=mystyle}

\begin{document}
	\noindent
	\begin{minipage}{1.7cm}
		{\includegraphics[trim=34.1cm 91.3cm 46.4cm 21cm,clip,width=.70\textwidth]{Marca_UFG.pdf}}
	\end{minipage}	
	\begin{minipage}{.8\textwidth}
		\center
		\hspace*{\fill}\\
		{\large Federal University of Goias} \\
		{\large Institute of Informatics} \\
		{\Large INF0417 -- Computer vision}
	\end{minipage}
	\begin{minipage}{2.2cm}
		{\includegraphics[trim=0cm 0cm 0cm 0cm,clip,width=.70\textwidth]{Marca_INF.pdf}} % marca INF
	\end{minipage}
%	\vspace{-1em}

	\begin{minipage}{\textwidth}
		\maketitle
	\end{minipage}

\section{Introduction}
During this tutorial, the goal is to become familiar with Python and the NumPy library.
You should also get a better feeling for how images are represented as matrices as well as the connection between mathematical expressions and the Python code to implement them.
We additionally provide a guide for how to get started with using an IDE for python, which will make your life much easier during the projects.

\subsection{Exercises}
\begin{enumerate}
	\item Make sure that you’ve downloaded the notebook for this lesson at SIGAA:
	\begin{center} \vspace{-1em}
		\texttt{INF0417\_python-tutorial.ipynb}
	\end{center}
	\item Start a Jupyter notebook from a terminal:
	\begin{center} \vspace{-1em}
		\texttt{jupyter notebook}
	\end{center}
	\item Find the correct notebook and open it.
	\item Follow the instructions in there.
\end{enumerate}

\section{Software tools}
This section covers unit tests, and development tools.
None of this is strictly part of the course, but they are very useful.

\subsection{Development tools}
There exist a large number of different development tools for Python, ranging from a text-editor like \texttt{atom} that can make it easy to get started, to full IDEs, such as
\texttt{pycharm}.
Getting started using an IDE such as \texttt{pycharm might} require some time to get familiar with the program, but it will usually save time in the long run due to the very useful auto-complete functions and integrated debugging functionality.

For the projects we strongly recommend that you use \texttt{pycharm} or \texttt{vscode}.

\begin{itemize}
	\item \textbf{Task 1:} Open \texttt{vscode}. Try creating a python file and running it.
	\item \textbf{Task 2:} Open \texttt{pycharm}. Try creating a python file and running it.
\end{itemize}

Both these IDEs will provide you with a good development environment.
Choose whichever you feel more comfortable with.

\subsection{Unit tests}
When writing programs larger than a few lines it is often useful to test the intermediate steps of the program in isolation.
There are (at least) two reasons for this:
\begin{enumerate}
	\item It helps you define what each part of the program should actually do.
	\item It makes sure that it actually does what you intended in isolation from the other parts.
\end{enumerate}


For larger projects involving multiple people, \textbf{unit tests} can allow you to test each component in isolation for some simple cases.
A unit-test is a function that calls another function you want to test with the arguments that should produce some known output, for example:

\lstinputlisting[language=Python, firstline=1, lastline=7]{INF0417\_lab-guidelines.py}

This function tests that the addition operation behaves as expected.
\\
Often you would write multiple tests so that you can be sure that edge cases are handled correctly.
A more complicated example for a function that projects a point in an image:

\lstinputlisting[language=Python, firstline=9, lastline=12]{INF0417\_lab-guidelines.py}

With the tests being:

\lstinputlisting[language=Python, firstline=14, lastline=24]{INF0417\_lab-guidelines.py}

Try to implement the missing test cases.
\begin{itemize}
	\item \textbf{Hint:} Use the same camera matrix in all tests, but different points.
\end{itemize}

Note that the point of this exercise is as much to implement the tests as it is to think about what the result should be in each case.

\printbibliography

\end{document}